\scanpage{75}
\mspage{II.23}
% archive sheet number: 74
% line 1
3. \uline{Réduction au cas de la dimension} 1.

% lines 2–6
3.1. \deleted{\unsure{Supp}} Soit d'abord $\unsure{k|h}$ un \deleted{\unsure{groupe}} corps quelconque, $X'$ un schéma \deleted{géométriquement irréductible} \inserted{\deleted{\unsure{de}} projectif \inserted{géom} normal, géom connexe,} sur $k$, $X$ un ouvert dense de $X$, \deleted{On se propose d'} \inserted{$\unsure{\xi|\eta}$ un pt géom. de $X'$, $\bar k$ la clôture séparable de $k$ dans $\kappa(\xi)$, $\bar X',\bar X$ déduits de $X',X$ par $k\unsure{-|\longrightarrow}\bar k$.} On veut étudier $\pi_1(\bar X^{\unsure{\sharp|*|\prime}},\xi)$, et \deleted{\unsure{de}} le majorer en fonction du groupe fondamental d'une \uline{courbe} lisse sur un corps convenable.
% lines 7–11
Comme on veut tenir compte des opérations, possibles de $\pi_1(S,\xi)$ ($S=\operatorname{Spec}k$) sur $\pi_1(\bar X^\sharp,\xi)$ resp sur sa partie abélienne, on va considérer aussi la suite exacte
% line 12
\noindent (3.1.1)
\diagram{\begin{tikzcd}[ampersand replacement=\&]
1 \arrow[r,dash] \& \pi_1(\bar X,\xi) \arrow[r,dash] \& \pi_1(X,\xi) \arrow[r,dash] \& \pi_1(S,\xi) \arrow[r,dash] \& 1 .
\end{tikzcd}}

% lines 13–19
\deleted{Soit $T$ le schéma qui} Nous supposons \deleted{que} $X'$ immergé dans un $\unsure{\mathbb P^r_k|P^r_k}$. \deleted{$\unsure{P^r_k}$ \unsure{Soit}} Soit $n=\dim X'$, et soit $T$ le grassmannien des sous-espaces projectifs de codimension $n-1$ \deleted{\unsure{de}} dans $P$, \inserted{(EGA I, $2^{\text{ième}}$ édition\unsure{\ldots})} $t$ le point générique de $T$, $K=k(t)$, qui est une extension \deleted{pure} \inserted{de t.f.} de $k$. \deleted{Alors} Soit $L_t$ le sous-espace \inserted{projectif} de codim $n-1$ \inserted{``générique''} de $P_K$ qui correspond à $t\in T_K$, et considérons $C'=L\cap X'_K$, \deleted{qui est} Il est connu que $\dim C'=1$, que $C'$ est géom. normal et géom connexe. Soit $C$ l'image inverse de $X_K$ dans $C'$.
% lines left of 18–21
\diagram{\begin{tikzcd}[ampersand replacement=\&]
C \arrow[d] \& C'=L_t\cap X'_K \arrow[l] \arrow[d,dash] \\
X_K \arrow[r,phantom,"\subset"] \& X'_K
\end{tikzcd}}
